% answers_short.tex —— 2009 年·普通高等学校招生考试·天津卷·理科 · 简短答案速查
% 由 tools/make_short_answers.py 自 answers.tex 的 \gksolhead/\gkanswer 提取，
% 并与 2009/2009_数学真题_天津理_答案要点.md 逐题比对后生成，勿手改（改源 MD 后重跑）。

\begin{gknotice}
\textbf{使用说明}：本节只给最终答案便于快速核对，完整推导见解析版。答案为编写者独立推导，\textbf{非官方标准答案}。
\end{gknotice}

\gkpart{选择题}
\begin{center}\begin{tabular}{ccccccccccc}
\toprule
题号 & 1 & 2 & 3 & 4 & 5 & 6 & 7 & 8 & 9 & 10 \\
\midrule
答案 & D & B & D & D & C & B & A & C & A & C \\
\bottomrule
\end{tabular}\end{center}

\gkpart{填空题}
\begin{itemize}[leftmargin=2.6em]
  \item[\textbf{11.}] $40$。
  \item[\textbf{12.}] $\sqrt3$。
  \item[\textbf{13.}] $\frac{3\sqrt{10}}{5}$。
  \item[\textbf{14.}] $1$。
  \item[\textbf{15.}] $\sqrt3$。
  \item[\textbf{16.}] $324$。
\end{itemize}

\gkpart{解答题}
\begin{itemize}[leftmargin=2.6em]
  \item[\textbf{17.}] （1）$AB=2\sqrt5$；（2）$\frac{\sqrt2}{10}$。
  \item[\textbf{18.}] （1）分布列见下，$EX=\frac9{10}$；（2）$\frac{31}{120}$。
  \item[\textbf{19.}] （1）$60^\circ$；（2）证明见解析；（3）$\frac{\sqrt3}{3}$。
  \item[\textbf{20.}] （1）$3\mathrm e$；（2）见下。
  \item[\textbf{21.}] （1）$\frac{\sqrt3}{3}$；（2）$\pm\frac{\sqrt2}{3}$；（3）$\pm\frac{2\sqrt2}{5}$。
  \item[\textbf{22.}] （1）$55$；（2）（3）证明见解析。
\end{itemize}
